%!TEX program = xelatex
\documentclass[a4paper,12pt]{article}
\usepackage{fontspec}\defaultfontfeatures{Ligatures=TeX}
\usepackage[a4paper,vmargin={3cm,3.5cm},hmargin={3cm,3cm}]{geometry}
%-----------------------------------------------------------------------------%
\usepackage{settings}
%-----------------------------------------------------------------------------%
%%% Title %%%
\title{G6411 Recitation Notes: \glsentryshort{ols} Large Sample Properties}
\author{Jesse Chieh Chen\\\href{mailto:cc5229@columbia.edu}{cc5229@columbia.edu}}
\date{\today}
%-----------------------------------------------------------------------------%

\begin{document}

\appendix
\section{Modes of Convergence}

Let $\{Z_n\}_{n=1}^{\infty}$ be a sequence of scalar random variables and let $Z$ be a scalar random variable.
A constant is also allowed as the limit.

\begin{definition}[Convergence in Probability]
	We say that $Z_n$ \emph{converges in probability} to $Z$, written $Z_n\pto Z$, if
	for every $\varepsilon>0$,
	\begin{align}\label{eq-probability-convergence}
		\Pr(|Z_n-Z|>\varepsilon)\to0
		\qquad\text{as }n\to\infty.
	\end{align}
\end{definition}

\begin{definition}[Convergence in Mean Square]
	We say that $Z_n$ \emph{converges in mean square} to $Z$, written $Z_n\msto Z$, if
	\begin{align}\label{eq-mean-square-convergence}
		\E[(Z_n-Z)^2]\to0,
	\end{align}
	provided that the variables have finite second moments.
\end{definition}

\begin{proposition}\label{prop-ms-implies-p}
	If $Z_n\msto Z$, then $Z_n\pto Z$.
\end{proposition}
\begin{proof}
	For any $\varepsilon>0$, Markov's inequality applied to the nonnegative random variable $(Z_n-Z)^2$ gives
	\begin{align*}
		\Pr(|Z_n-Z|>\varepsilon)
		=\Pr((Z_n-Z)^2>\varepsilon^2)
		\le\frac{\E[(Z_n-Z)^2]}{\varepsilon^2}
		\to0.
	\end{align*}
	This is the definition of convergence in probability.
\end{proof}

\begin{remark}[Bias-Variance Decomposition]
	For a constant target $c$, the same result follows from the bias-variance decomposition:
	\begin{align}\label{eq-mse-decomposition}
		\E[(Z_n-c)^2]=\var(Z_n)+(\E(Z_n)-c)^2.
	\end{align}
	If variance and squared bias both vanish, the \gls{mse} tends to zero.
	Mean square convergence then implies convergence in probability by \autoref{prop-ms-implies-p}.
\end{remark}

\begin{example}[Converse does not hold]\label{ex-p-not-ms}
	Let $\{Z_n\}$ be a sequence of random variables with $\supp(Z_n)=\{0,n\}$ $\forall n$.
	Let $Z_n=n$ with probability $1/n$ and $Z_n=0$ otherwise.
	For every fixed $\varepsilon>0$ and sufficiently large $n$,
	\begin{align*}
		\Pr(|Z_n|>\varepsilon)=n\inv\to0
	\end{align*}
	Thus $Z_n\pto0$.
	However, since $\E(Z_n^2)=n$, the sequence does not converge to zero in mean square.
\end{example}

\begin{definition}[Convergence in Distribution]\label{def-distribution-convergence}
	Let $F_{Z_n}$ and $F_Z$ be the \gls{cdf} of $Z_n$ and $Z$ respectively.
	We say that $Z_n$ \emph{converges in distribution} to $Z$, written $Z_n\dto Z$, if
	\begin{align}
		F_{Z_n}(z)\to F_Z(z)
		\qquad\text{for all }z\in\reals\text{ at which }F_Z(z)\text{ is continuous}.
	\end{align}
\end{definition}

\begin{remark}
	One technical difference between convergence in probability and convergence in distribution is that
	the latter does not require the random variable to be defined on the same probability space,
	as we only require the \glspl{cdf} to converge at continuous points.
\end{remark}

\begin{example}
	Consider the sequence of random variables $Z_n$ with \glspl{cdf}
	\begin{align}
		F_{Z_n}(z) =
		\begin{cases}
			0, & z < -\frac1n \\
			\frac12 + \frac{n}{2} z, & z \in \left[-\frac1n,\frac1n\right] \\
			1, & z > \frac1n
		\end{cases}
	\end{align}
	It is easy to see that $F_{Z_n}(z)\to F_Z(z)$ for all $z\neq0$, where $F_Z$ is the \gls{cdf} of the constant random variable $Z=0$.
	Note that $F_{Z_n}(0)=1/2\neq F_Z(0)=1$.
\end{example}

\begin{proposition}
	If $Z_n\pto Z$, then $Z_n\dto Z$.
\end{proposition}
\begin{proof}
	Fix a point $z$ at which $F_Z$ is continuous and let $\delta>0$.
	If $|Z_n-Z|\le\delta$, then $Z\le z-\delta$ implies $Z_n\le z$,
	and $Z_n\le z$ implies $Z\le z+\delta$.
	Consequently,
	\begin{align*}
		F_Z(z-\delta)-\Pr(|Z_n-Z|>\delta)
		\le F_{Z_n}(z)
		\le F_Z(z+\delta)+\Pr(|Z_n-Z|>\delta).
	\end{align*}
	Since $Z_n\pto Z$, the probability terms tend to zero, giving
	\begin{align*}
		F_Z(z-\delta)
		\le\liminf_{n\to\infty}F_{Z_n}(z)
		\le\limsup_{n\to\infty}F_{Z_n}(z)
		\le F_Z(z+\delta).
	\end{align*}
	Letting $\delta\to0$, continuity at $z$ makes both outer bounds converge to $F_Z(z)$.
	Thus $F_{Z_n}(z)\to F_Z(z)$ at every continuity point of $F_Z$.
\end{proof}

\begin{proposition}\label{prop-d-to-p}
	If $Z_n\dto c$ for a constant $c$, then $Z_n\pto c$.
\end{proposition}
\begin{proof}
	The \gls{cdf} of the constant $c$ is zero below $c$ and one at or above $c$.
	For any $\varepsilon>0$, both $c-\varepsilon$ and $c+\varepsilon$ are continuity points of this \gls{cdf}.
	Convergence in distribution therefore gives
	$F_{Z_n}(c-\varepsilon)\to0$
	and
	$F_{Z_n}(c+\varepsilon)\to1$.
	Hence,
	\begin{align*}
		\Pr(|Z_n-c|>\varepsilon)
		&=\Pr(Z_n<c-\varepsilon)+\Pr(Z_n>c+\varepsilon)\\
		&\le F_{Z_n}(c-\varepsilon)+1-F_{Z_n}(c+\varepsilon)\to0.
		\qedhere
	\end{align*}
\end{proof}

\begin{remark}[Chain of Implications]
	Let $Z_n$ be a sequence of random variables and $Z$ be a random variable or constant.
	Then we have the following chain of implications:
	\begin{align}
		Z_n\msto Z
		\implies Z_n\pto Z
		\implies Z_n\dto Z.
	\end{align}
	The converse does not hold in general.
	\autoref{ex-p-not-ms} shows that the first implication is strict;
	And unless $Z$ is a constant, convergence in distribution does not imply convergence in probability.
\end{remark}

\section{Useful Theorems}

\begin{lemma}[Concentration around the mean]\label{lem-vanishing-variance}
	Let $\{Z_n\}$ be a sequence of random variables with finite variance
	for every $n$ and $\var(Z_n)\to0$.
	Then
	\begin{align}
		Z_n-\E(Z_n)\pto0.
	\end{align}
	If also $\E(Z_n)\to c$ for a constant $c$, then $Z_n\pto c$.
\end{lemma}
\begin{proof}
	Chebyshev's inequality gives, for every $\varepsilon>0$,
	\begin{align*}
		\Pr(|Z_n-\E(Z_n)|>\varepsilon)
		\le\frac{\var(Z_n)}{\varepsilon^2}\to0.
	\end{align*}
	If $\E(Z_n)\to c$, then for sufficiently large $n$,
	$|\E(Z_n)-c|<\varepsilon/2$.
	Hence, we have
	\begin{align*}
		\Pr(|Z_n-c|>\varepsilon)
		&\le \Pr(|Z_n-\E(Z_n)|+|\E(Z_n)-c|>\varepsilon) \\
		&\le \Pr(|Z_n-\E(Z_n)|>\varepsilon/2)
		\le \frac{4\var(Z_n)}{\varepsilon^2}\to0.
		\qedhere
	\end{align*}
\end{proof}

\begin{theorem}[\gls{wlln}]\label{thm-wlln}
	Let $\{X_i\}_{i=1}^{n}$ be \gls{iid} with mean $\mu$ and finite variance $\sigma^2$.
	Then the sample mean $\bar X_n=n\inv\sum_{i=1}^nX_i$ is consistent for $\mu$.
\end{theorem}
\begin{proof}
	Apply \autoref{lem-vanishing-variance} with $Z_n=\bar X_n$ and $c=\mu$.
\end{proof}

\begin{remark}
	The finite second moment assumption in \autoref{thm-wlln} can be dropped.
	In fact, \gls{wlln} holds for \gls{iid} random variables with finite first moment.
	That is, if $\{X_i\}_{i=1}^{n}$ is \gls{iid} with $\E(X_i)=\mu$ finite, then $\bar X_n\pto\mu$.
	However, the proof of this stronger result is more involved.
\end{remark}

\begin{theorem}[Continuous Mapping]\label{thm-continuous-mapping}
	Let $Z_n\in\reals^k$ be random vectors and $Z\in\reals^k$ be a random vector.
	Let $g$ be a continuous function.
	Then we have
	\begin{align}
		Z_n\pto Z  &\implies g(Z_n)\pto g(Z), \\
		Z_n\dto Z  &\implies g(Z_n)\dto g(Z).
	\end{align}
\end{theorem}
\begin{proof}
	Omitted.
\end{proof}

\begin{example}
	If $(Z_{1n},Z_{2n})\T\pto (Z_{1,0},Z_{2,0})\T\in\reals^2$.
	By \autoref{thm-continuous-mapping}, we have
	\begin{align}
		Z_{1n}+Z_{2n}&\pto Z_{1,0}+Z_{2,0}, \\
		Z_{1n}Z_{2n}&\pto Z_{1,0}Z_{2,0}, \\
		Z_{1n}/Z_{2n}&\pto Z_{1,0}/Z_{2,0}\quad(Z_{2,0}\neq0)
	\end{align}
	by applying the continuous mapping theorem to the functions
	$g(x,y)=x+y$, $g(x,y)=xy$, and $g(x,y)=x/y$ respectively to $(Z_{1n},Z_{2n})$.
\end{example}

\begin{theorem}[Slutsky's Theorem]\label{thm-slutsky}
	Let $Z_n$ and $W_n$ be sequences of random variables.
	If $Z_n\dto Z$ and $W_n\pto c$ for a constant $c$, then
	\begin{align}
		Z_n+W_n\dto Z+c,
		\qquad
		Z_nW_n\dto cZ,
		\qquad
		Z_n/W_n\dto Z/c\quad(c\neq0).
	\end{align}
\end{theorem}
\begin{proof}[Sketch of Proof]
	Convergence in probability implies $W_n\dto c$.
	One can show that $(Z_n,W_n)\T\dto(Z,c)\T$.
	Then apply \autoref{thm-continuous-mapping} to the functions $g(x,y)=x+y$, $g(x,y)=xy$, and $g(x,y)=x/y$ respectively to $(Z_n,W_n)\T$.
\end{proof}

\begin{theorem}[Lindeberg-Lévy \gls{clt}]\label{thm-clt}
	Let $\{X_i\}_{i=1}^{n}$ be \gls{iid} with mean $\mu$ and finite variance $\sigma^2$.
	Then
	\begin{align}
		\sqrt{n}(\bar X_n-\mu)\dto\normaldistribution(0,\sigma^2).
	\end{align}
\end{theorem}
\begin{proof}
	Omitted.
\end{proof}

\begin{theorem}[Delta Method]\label{thm-delta-method}
	Let $Z_n\in\reals^k$ be a sequence of random vectors and $Z\in\reals^k$ be a fixed vector.
	Let $g:\reals^k\to\reals^m$ be continuously differentiable in a neighborhood of $Z$, with Jacobian matrix $G(Z)\in\reals^{m\times k}$.
	If $\sqrt{n}(Z_n-Z)\dto Z_0$, then
	\begin{align}
		\sqrt{n}(g(Z_n)-g(Z))\dto G(Z)Z_0.
	\end{align}
\end{theorem}
\begin{proof}
	Since $\sqrt{n}(Z_n-Z)=O_p(1)$, we have $Z_n-Z=O_p(n^{-1/2})$ and hence $Z_n\pto Z$.
	Since $Z_n\pto Z$, it suffices to work in the neighborhood where $g$ is continuously differentiable.
	The integral form of the mean value theorem gives the matrix identity
	\begin{align*}
		g(Z_n)-g(Z)=H_n(Z_n-Z)
		&\quad\text{where}\quad
		H_n\coloneqq h(Z_n)\coloneqq \int_0^1G\big(Z+t(Z_n-Z)\big)\diff{t}.
	\end{align*}
	Continuity of $G$ makes $h$ continuous at $Z$ with $h(Z)=G(Z)$.
	Since $Z_n\pto Z$, \nameref{thm-continuous-mapping} gives $h(Z_n)\pto h(Z)$, i.e., $H_n-G(Z)=o_p(1)$.
	Hence, we have
	\begin{align*}
		\sqrt{n}(g(Z_n)-g(Z))
		&= \sqrt{n} H_n(Z_n-Z) \\
		&= G(Z)\sqrt{n}(Z_n-Z) + \underbrace{\sqrt{n} \underbrace{(H_n-G(Z))}_{o_p(1)} \underbrace{(Z_n-Z)}_{O_p(n^{-1/2})}}_{o_p(1)}.
	\end{align*}
	The first term converges in distribution to $G(Z)Z_0$ by the continuous mapping theorem,
	while the second converges in probability to zero.
	\nameref{thm-slutsky} gives the stated result.
\end{proof}

\end{document}
